% PhysXWrapper Technical Documentation
% Author: PhysX Wrapper Project
% Date: 2025-11-06

\documentclass[12pt,a4paper]{report}

% 包导入
\usepackage[UTF8]{ctex}  % 中文支持
\usepackage{geometry}
\usepackage{graphicx}
\usepackage{tikz}
\usepackage{listings}
\usepackage{xcolor}
\usepackage{hyperref}
\usepackage{tabularx}
\usepackage{longtable}
\usepackage{amsmath}
\usepackage{amssymb}
\usepackage{fancyhdr}
\usepackage{tocloft}

% TikZ库
\usetikzlibrary{shapes.geometric, arrows.meta, positioning, calc, shadows, decorations.pathreplacing}

% 页面设置
\geometry{left=3cm,right=2.5cm,top=2.5cm,bottom=2.5cm}

% 代码高亮设置
\lstdefinestyle{cpp}{
    language=C++,
    basicstyle=\ttfamily\small,
    keywordstyle=\color{blue}\bfseries,
    commentstyle=\color{green!60!black}\itshape,
    stringstyle=\color{red},
    numbers=left,
    numberstyle=\tiny\color{gray},
    stepnumber=1,
    numbersep=8pt,
    backgroundcolor=\color{gray!10},
    showspaces=false,
    showstringspaces=false,
    showtabs=false,
    frame=single,
    rulecolor=\color{black!30},
    tabsize=4,
    captionpos=b,
    breaklines=true,
    breakatwhitespace=false,
    escapeinside={\%*}{*)},
    xleftmargin=2em,
    xrightmargin=2em,
    aboveskip=1em,
    belowskip=1em
}

\lstset{style=cpp}

% 超链接设置
\hypersetup{
    colorlinks=true,
    linkcolor=blue,
    filecolor=magenta,
    urlcolor=cyan,
    citecolor=green,
    pdftitle={PhysXWrapper Technical Documentation},
    pdfauthor={PhysX Wrapper Project},
    pdfsubject={Physics Engine Wrapper},
    pdfkeywords={PhysX, Physics, C++, API}
}

% 页眉页脚
\pagestyle{fancy}
\fancyhf{}
\fancyhead[L]{\leftmark}
\fancyhead[R]{PhysXWrapper}
\fancyfoot[C]{\thepage}
\renewcommand{\headrulewidth}{0.4pt}
\renewcommand{\footrulewidth}{0.4pt}

% 标题页设置
\title{
    \vspace{2cm}
    \Huge \textbf{PhysXWrapper} \\
    \vspace{0.5cm}
    \LARGE 技术文档与API参考手册 \\
    \vspace{0.3cm}
    \Large Technical Documentation \& API Reference \\
    \vspace{2cm}
}
\author{
    PhysX Wrapper Project Team \\
    \vspace{0.5cm}
    \texttt{基于 NVIDIA PhysX 5.6.1}
}
\date{\today}

\begin{document}

% 标题页
\maketitle
\thispagestyle{empty}

% 摘要
\newpage
\begin{abstract}
\noindent
PhysXWrapper是对NVIDIA PhysX物理引擎的高级C++封装库，旨在简化PhysX的使用，提供更加友好和易用的接口。本文档详细记录了PhysXWrapper的架构设计、API参考、示例程序库、代码移植对比以及性能分析。

\vspace{1em}
\noindent
\textbf{主要特性：}
\begin{itemize}
    \item 完整实现72个PhysX官方Snippets (100\%覆盖率)
    \item 模块化设计，支持刚体、软体、流体、车辆等
    \item RAII资源管理，自动生命周期控制
    \item 详细的中文注释和理论背景
    \item 完善的单元测试和集成测试
\end{itemize}

\vspace{1em}
\noindent
\textbf{文档版本：} 1.0 \\
\textbf{最后更新：} \today \\
\textbf{PhysX版本：} 5.6.1 \\
\textbf{C++标准：} C++17
\end{abstract}

% 目录
\newpage
\tableofcontents

% 图表目录
\newpage
\listoffigures
\listoftables

% 正文章节
\newpage
\input{chapters/01_introduction}
\input{chapters/02_api_reference}
\input{chapters/03_example_library}
\input{chapters/04_code_comparison}
\input{chapters/05_architecture}
\input{chapters/06_performance}

% 附录
\appendix
\input{chapters/appendix_build}
\input{chapters/appendix_troubleshooting}

% 参考文献
\begin{thebibliography}{99}
\bibitem{physx} NVIDIA Corporation. \textit{PhysX SDK Documentation}. 2024.
\bibitem{physx_guide} NVIDIA Corporation. \textit{PhysX Programming Guide}. 2024.
\bibitem{snippets} NVIDIA Corporation. \textit{PhysX Snippets Collection}. 2024.
\bibitem{cpp17} ISO/IEC 14882:2017. \textit{Programming Languages - C++}. 2017.
\end{thebibliography}

\end{document}
